\documentclass[10pt]{article}

\usepackage[utf8]{inputenc}

\usepackage[T1]{fontenc}

\usepackage{amsmath}

\usepackage{amsfonts}

\usepackage{amssymb}

\usepackage[version=4]{mhchem}

\usepackage{stmaryrd}

\usepackage{graphicx}

\usepackage[export]{adjustbox}

\graphicspath{ {./images/} }



\begin{document}

н о н е п р е р ы в н м (относительно равномерности $\mathfrak{A})$, если для любой точки $x \in X$ и любого $V \in \mathfrak{A}$ существует окрестность $O_{x} \ni x$ такая, что $\left(f(x), f\left(x^{\prime}\right)\right) \in V$ при $x^{\prime} \in O_{x}$ и $f \in F$. Имеет место следующее обобщениө классич. теоремы Асколи: пусть $X$ есть $k$-пространство, $(Y, \mathfrak{A})$ - равномерное пространство и $\boldsymbol{Y}^{X}-$ пространство нешрерывных отображений $X$ в $Y$ с компактно открытой топологией. Для того чтобы замкнутое подмножество $F \subset Y X$ было бикомпактным, необходимо и достаточно, чтобы $F$ было равностепенно непрерывно относительно равномерности $\mathfrak{Y}$ и все множества $\{f(x)$ : $f \in F\}, \quad x \in X$, имели бикомпактные замыкания в $Y$ ( $k$-пространства - это хаусдорфовы пространства, являющиеся факторным образом локально компактных пространств; класс $k$-пространств содержит все хаусдорфовы пространства с первой аксиомой счетности и все хаусдорфовы локально бикомпактные пространства).



Топология метризуемого Р. п. паракомпактна в силу теоремы Стоуна. Однако решается отрицательно проблема Исбелла о равномерной паракомпактности метризуемых Р. п. Построен пример метризуемого Р. П., имеющего равномерное покрытие.



В теории размерности Р. П. основными являются равномерные размерностные инварианты $\delta d$ и $\Delta d$, определяемые аналогично топологич. размерности $\mathrm{dim}(\delta \mathrm{d}-$ при помощи конечных равномерных покрытий, а $\Delta d-$ при помощи всех равномерных покрытий), и равномерная индуктивная размерность $\delta$ Ind. Размерность $\delta$ Ind определяется по аналогии с большой индуктивной размерностью Ind индукцией по размерности близостных перегородок между далекими (в смысле близости, порожденной равномерностью) множествами. При этом множество $H$ наз. б ли и о с т н о й п е р е г о р о дк о й между $A$ и $B$ (где $A \delta B$, если для любой $\delta$-окрестности $U$ множества $H$ такой, что $U \cap(A \cup B) \neq \varnothing$, имеет место $X \backslash U=A^{\prime} \cup B^{\prime}$, где $A^{\prime} \bar{\delta} B^{\prime}, A \subset A^{\prime}, B \subset B^{\prime}$ ( $U$ наз. $\delta$-окрестностью $H$, если $H \bar{\delta}(X \backslash U)$. Таким образом, размерность $\delta$ Ind (равно как и $\delta d$ ) является не только равномерным, но и близостным инвариантом. Размерность $\delta d$ P. п. $(X, U)$ совпадает с обычной размервостью $\operatorname{dim}$ расширения Самюəля, построенного по эквивалентной $\mathfrak{A}$ прекомпактной равномерности. Если размерность $\Delta \mathrm{d} X$ конечна, то $\Delta \mathrm{d} X=\delta \mathrm{d} X$. Однако может быть $\delta \mathrm{d} X=0$, a $\Delta \mathrm{d} X=\infty$. Для метризуемого Р. п. всегда $\delta \mathrm{d} X \leqslant \delta$ Ind $X=\Delta \mathrm{d} X$ (и если $\Delta \mathrm{d} X<\infty$, то $\delta \mathrm{d} X=\delta \operatorname{Ind} X=\Delta \mathrm{d} X)$. Равенства $\delta \mathrm{d} X=0$ и $\delta \operatorname{Ind} X=0$ эквивалентны для любого Р. п. Если Р. п. метризуемо, то эквивалентны также равенства $\delta \mathrm{d} X=0$ и $\Delta \mathrm{d} X=0$. Если P. п. $X^{\prime}$ является всюду плотным подмножеством P. п. $X$, тo $\delta$ Ind $X^{\prime} \geqslant \delta \operatorname{Ind} X$. Bcerдa $\delta I$ Ind $X \leqslant \operatorname{Ind} X$. Для размерности $\delta d$ имеет место аналог теоремы о перегородках.



Различные обобщения Р. п. получаются путем ослабления аксиом равномерности. Так, в аксиоматике квазиравномерности (см. [8]) исключена аксиома симметрии. При определении обобщенной равномерности (см. [10]) ( $f$-равномерности) вместо равномерных покрытий рассматриваются равномерные семейства подмножеств $X$, к-рые, вообще говоря, не являются покрытиями (тела этих семейств оказываются всюду плотными в $X$ в топологии, порожденной $f$-равномерностью). Одно из обобщений равномерности - так наз. $\theta-p$ а в н ом е р н о с т в - связано с наличием топологии на Р. п. Она определяется семействами $\theta$-покрытий хаусдорфова пространства; $\theta$-покрытиями наз. система $\eta$ канонических открытых множеств $X$, удовлетворяющих следующему условию: для любой точки $x \in X$ существуют $V_{1}, \ldots, V_{n} \in \eta$ такие, что $x \in \operatorname{int} \bigcup_{i=1}^{n}\left[V_{i}\right]$.



Jum.: [1] W e i l A., Sur les espaces à structure uniforme et sur le topologie générale, P., 1937; [2] Б у р б а к и Н., Общая топология. Основныө структуры, пер. с франц., 2 изд., М., 1968; [3] е г о ж е, Общая топология. Использование вещественных чисел в общей топологии. Функциональные пространства. Сводка результатов. Словарь, пер. с франц., M , 1975, [4]



\includegraphics[max width=\textwidth, center]{2023_07_08_50d4a69b290bddc79ea0g-1}

т p я г и н J. С., Непрерывные группы, 3 изд., M.; $1973 ;[6]$

I s b e l l J., Uniform spaces, Providence, $1964 ;[7]$ a m ue l P., "Trans. Amer. Math. Soc.", 1948, v. 64, p. $100-32 ;$ [8] C s á śz á r A., Foundations of general topology, Oxford-[etc.], $1963 ;[9] \Phi$ е д о р ч у к B. В., «Докл. АН СССР», 1970, т. 192 № 6; c. $1228-30 ;[10]$ K u 1 p a W., "Colloq. Math.", 1972, v. 25



\begin{center}

\includegraphics[max width=\textwidth]{2023_07_08_50d4a69b290bddc79ea0g-1(1)}

\end{center}



РАВНОМЕРНОЕ РАСПРЕДЕЛЕНИЕ - общее название класса распределений вероятностей, возникающөго при распространөнии идеи «равновозможности исходов» на нөпрерывный случай. Подобно нормальному распределению Р. р. появляөтся в теории вероятностей как точное распределение в одних задачах и как предельное - в других.



P. p. н о трезке чис ло во й пря мо й. (пря я о у гольно в P. р. на каком-либо отрезке $[a, b], a<b$, - это распределений вероятностей, имеющее плотность



\[

p(x)=\left\{\begin{array}{cl}

\frac{1}{b-a}, & x \in[a, b], \\

0, & x \notin[a, b] .

\end{array}\right.

\]



Понятие Р. p. на $[a, b]$ соответствует представлению о случайном выборе точки на әтом отрезке «наудачу». Математич. ожидание и дисперсия Р. р. равны, соответственно, $(b+a) / 2$ и $(b-a)^{2} / 12$. Функция распределения задается формулой



\[

F(x)=\left\{\begin{array}{cl}

0, & x \leqslant a \\

\frac{x-a}{b-a}, & a<x \leqslant b, \\

1, & x>b,

\end{array}\right.

\]



а характеристич. Функция - формулой



\[

\varphi(t)=\frac{1}{\imath t(b-a)}\left(e^{i t b}-e^{i t a}\right) .

\]



Случайную величину с Р. p. на $[0,1]$ можно построить, исходя из последовательности независимых случайных величин $X_{1}, X_{2}, \ldots$, принимающих значения 0 и 1 с вероятностями $1 / 2$, полагая



\[

X=\sum_{n=1}^{\infty} X_{n} 2^{-n}

\]



$\left(X_{n}\right.$ являются цифрами в двоичном разложении $\left.X\right)$. Случайное число $X$ имеет P. р. на отрезке $[0,1]$. Этот факт имеет важные статистич. приложения, см., напр., Случайные и псевдослучайные числа.



Если независимые случайные величины $X_{1}$ и $X_{2}$ имеют P. p. на $[0,1]$, то их сумма $X_{1}+X_{2}$ имеет так наз. треугольно е ра спределени е на $[0,2]$ с плотностью $u_{2}(x)=1-|1-x|$ для $x \in[0,2]$ и $u_{2}(x)=0$ для $x \notin[0,2]$. Сумма трех независимых случайных величин c P. p. на $[0,1]$ имеет распределение на $[0,3]$ с плотностью



\[

u_{3}(x)= \begin{cases}x^{2} / 2, & 0 \leqslant x<1, \\ {\left[x^{2}-3(x-1)^{2}\right] / 2,} & 1 \leqslant x<2, \\ {\left[x^{2}-3(x-1)^{2}+3(x-2)^{2}\right] / 2,} & 2 \leqslant x<3, \\ 0, & x \notin[0,3] .\end{cases}

\]



В общем случае сумма $X_{1}+X_{2}+\ldots+X_{n}$ независимых величин с P. p. на $[0,1]$ распределена с плотностью



\[

u_{n}(x)=\frac{1}{(n-1) !} \sum_{k=0}^{n}(-1)^{k} C_{n}^{k}(x-k)_{+}^{n-1}

\]



для $0 \leqslant x \leqslant n$ и $u_{n}(x)=0$ для $x \notin[0, n] ;$ здесь



\[

z_{+}= \begin{cases}z, & z>0 \\ 0, & z \leqslant 0\end{cases}

\]





\end{document}